\chapter{Einführungen}
\quellseite{10}

Im Auftrag des Ehrw. Vicepräses Schulz wurde P. Albert Brauer am Sonnt. n. Neujahr in seiner Gemeinde zu Fairfield, Minn., eingeführt von W. Marth.

Im Auftrag des Ehrw. Präses Becker wurde P. M. W. Fleckenstein am Sonnt. n. Neujahr in seiner Gemeinde zu Minden, Nebr., eingeführt von H. Hopmann.

Im Auftrag des Ehrw. Präses des Westlichen Districts wurde P. F. Möller am Sonnt. n. Neujahr in der Gemeinde zu Mora, Mo., unter Assistenz P. Wittrocks eingeführt von W. Matuschka.

Erhaltenem Auftrag gemäß wurde P. C. H. Polzin am Sonnt. n. Neujahr in der Parochie Otis, Westville und Porter, Ind., unter Assistenz der PP. Clausen und Schlechte eingeführt von A. Rump.

Im Auftrag des Ehrw. Präses Clöter wurde P. Martin Brüggemann am 1. Sonnt. n. Epiph. in seiner Gemeinde zu Deer Creek Township, Iowa, eingeführt von A. H. Deletzke.

Im Auftrag des Ehrw. Präses Walter wurde P. Martin Gallmeier am 1. Sonnt. n. Epiph. in seiner Gemeinde zu Cumberland, Md., eingeführt von S. Kliernericht.

Im Auftrag des Ehrw. Präses Clöter wurde P. A. Hemann am 1. Sonnt. n. Epiph. in seiner Gemeinde zu Atlantic, Iowa, eingeführt von R. Kolb.

Im Auftrag des Ehrw. Präses des Nebraska-Districts wurde P. F. Ruphoff am 2. Sonnt. n. Epiph. unter Assistenz der PP. Scholder, Eckhard und Schabacker in der Immanuels-Gemeinde bei Deshler, Nebr., eingeführt von S. Meyer.

Am 2. Sonnt. n. Epiph. wurde P. A. F. Scharfenberg in Lena und Richland Township, Ill., eingeführt von A. C. Landeck.

Im Auftrag des Ehrw. Präses Hafner wurde P. H. D. Wagner am 2. Sonnt. n. Epiph. in seinen Gemeinden zu Chepstow und Winklers Mills, Kanf., eingeführt von J. G. B. Keller.

Im Auftrag des Ehrw. Präses Riemann wurde P. A. Hänfschel am 2. Sonnt. n. Epiph. in seiner Gemeinde zu Columbus, O., unter Assistenz P. C. A. Brüggemanns eingeführt von H. W. Born.

Im Auftrag des Ehrw. Präses Clöter wurde P. D. Gurichte am 3. Sonnt. n. Epiph. in seiner Gemeinde bei Coon Rapids, Iowa, eingeführt von H. B. Säger.

Am 1. Sonnt. n. Epiph. wurde Lehrer P. W. Naßte in der Zion-Gemeinde zu Tobias, Nebr., eingeführt von Ch. Hartmann.

Am 2. Sonnt. n. Epiph. wurde Lehrer C. Schuhmacher als zweiter Lehrer an der Schule der St. Jacobi-Gemeinde zu Quincy, Ill., eingeführt von W. Hallerberg jun.

Am 3. Sonnt. n. Epiph. wurde Lehrer Leo H. Sippel, berufen an die Unterklasse der St. Johannes-Schule zu Vincennes, Ind., eingeführt von C. Kretzmann.
